
\section{诺维科夫环}
\label{sec:5-1}

\subsection{约化至万有诺维科夫环}
\label{subsec:5-1-1}

本书以万有诺维科夫环 $\Lambda_{\mathrm{nov}}$
及其子环 $\Lambda_{0,\mathrm{nov}}$
作为弗洛尔上同调的系数环，
并以 $R$ 表示万有诺维科夫环的底环。
第 2 章和其他若干文献使用不同类型的诺维科夫环研究弗洛尔上同调。
本节澄清它们之间的关系；
不过，本书不会用到关于其他类型诺维科夫环的这些结果。

先考虑单个拉格朗日子流形 $L$ 的情形。
回顾定义 2.4.17：在 $\pi_2(M,L)$ 中，
当且仅当 $E(\beta)=E(\beta')$ 且
$\mu_L(\beta)=\mu_L(\beta')$ 时，
定义 $\beta\sim\beta'$。
这里 $\mu_L$ 是马斯洛夫指标同态，
$E(\beta)=\int_\beta\omega$ 是辛面积。
相应等价类的集合记为
$\Pi(L)=\pi_2(M,L)/{\sim}$。
置
\[
\begin{aligned}
\Lambda(L)=
\biggl\{\sum_{\beta\in\Pi(L)}a_\beta[\beta]\;\bigg|\;&
a_\beta\in R,\ \text{且对每个 }C\in\mathbb R,\\
&\#\{\beta\mid a_\beta\ne0,\ E(\beta)<C\}<\infty
\biggr\}.
\end{aligned}
\]
第 2 章所述的弗洛尔上链复形 $CF(L,L)$
是 $\Lambda(L)$ 上的模。
定义同态 $I_L:\Lambda(L)\to\Lambda_{\mathrm{nov}}$：
\[
I_L\left(\sum_{\beta\in\Pi(L)}
a_\beta[\beta]\right)
=\sum_{\beta\in\Pi(L)}
a_\beta T^{E(\beta)}e^{\mu_L(\beta)/2}.
\]
下述引理由定义立即推出。

\begin{lemma}{5.1.1}\label{lem:5-1-1}
设 $\psi:M\to M'$ 是辛微分同胚，
且 $\psi(L)=L'$。
则映射
$\psi_*:\pi_2(M,L)\to\pi_2(M',L')$
典范地诱导同态
\[
\psi_*:\Pi(L)\to\Pi(L'),
\qquad
\psi_*:\Lambda(L)\to\Lambda(L').
\]
此外，有恒等式
$I_{L'}\circ\psi_*=I_L$。
\end{lemma}

现在回顾拉格朗日子流形对
$(L^{(0)},L^{(1)})$ 的诺维科夫环
$\Lambda(L^{(0)},L^{(1)};\ell_0)$
的定义。
第 2.2 节引入了从 $L^{(0)}$ 连到 $L^{(1)}$
的道路空间 $\Omega(L^{(0)},L^{(1)})$。
选取基道路 $\ell_0$，
并取其所在的连通分支
$\Omega(L^{(0)},L^{(1)};\ell_0)$。
定义 2.2.4 构造了该分支的覆盖空间
$\widetilde\Omega(L^{(0)},L^{(1)};\ell_0)$，
并以 $\Pi(L^{(0)},L^{(1)};\ell_0)$
表示这一覆盖的覆盖变换群。

回顾第 2.2.1 小节通过辛面积与马斯洛夫指标
分别构造的同态
\[
\begin{aligned}
E_{\ell_0}&:\Pi(L^{(0)},L^{(1)};\ell_0)
  \longrightarrow\mathbb R,\\
\mu_{\ell_0}&:\Pi(L^{(0)},L^{(1)};\ell_0)
  \longrightarrow\mathbb Z.
\end{aligned}
\]
推论 2.2.5 前的一段将它们简记为 $E$ 与 $\mu$；
这里把 $\ell_0$ 写在下标中，以表明对基道路选取的依赖。
若 $L^{(0)}$ 与 $L^{(1)}$ 均可定向，
则 $\operatorname{Image}(\mu_{\ell_0})\subset2\mathbb Z$。

把 $\Lambda^k(L^{(0)},L^{(1)};\ell_0)$
定义为形如
\[
\sum_{\substack{\beta\in\Pi(L^{(0)},L^{(1)};\ell_0)\\
                 \mu_{\ell_0}(\beta)=k}}
a_\beta[\beta]
\]
的形式和所成的集合，其中 $a_\beta\in R$，
且对每个常数 $C\in\mathbb R$，集合
\[
\{\beta\in\Pi(L^{(0)},L^{(1)};\ell_0)
\mid E_{\ell_0}(\beta)\leq C,\ a_\beta\ne0\}
\]
都是有限集。于是
\[
\Lambda(L^{(0)},L^{(1)};\ell_0)
=\bigoplus_k\Lambda^k(L^{(0)},L^{(1)};\ell_0).
\]
参见定义 2.2.6。
第 2.3、2.4 节定义的弗洛尔上链复形
$CF(L^{(1)},L^{(0)};\ell_0)$
是此诺维科夫环\liftIndex{I064}
$\Lambda(L^{(0)},L^{(1)};\ell_0)$ 上的模。
事实上，作为分次模，
$CF(L^{(1)},L^{(0)};\ell_0)$
还依赖于 $\lambda^0$ 的选取；
$\lambda^0$ 的定义见第 2.2.2 小节。
为强调这种依赖，下文写作
\[
CF(L^{(1)},L^{(0)};\ell_0)
=CF(L^{(1)},L^{(0)};\ell_0,\lambda^0).
\]
但诺维科夫环
$\Lambda(L^{(0)},L^{(1)};\ell_0)$
并不依赖于 $\lambda^0$ 的选取。

下面考察
$\Lambda(L^{(0)},L^{(1)};\ell_0)$
怎样依赖于 $\ell_0$。
按定义，它是群环
\[
G(\Pi(L^{(0)},L^{(1)};\ell_0))
\]
关于能量滤过的完备化，
因此只须考察
$\Pi(L^{(0)},L^{(1)};\ell_0)$
对 $\ell_0$ 的依赖。
又由该群的定义，
\[
\Pi(L^{(0)},L^{(1)};\ell_0)
=
\frac{\pi_1(\Omega(L^{(0)},L^{(1)});\ell_0)}
{\operatorname{Ker}E_{\ell_0}
 \cap\operatorname{Ker}\mu_{\ell_0}}.
\]
取一条道路
$\gamma:[0,1]\to\Omega(L^{(0)},L^{(1)})$，
连接 $\ell_0$ 与 $\ell_0'$。\translatornote{原书实体页 277
在上文群环以及本商式的道路空间中
都把 $(L^{(0)},L^{(1)})$ 倒写；
但基道路 $\ell_0$ 属于正序道路空间，
且第 2.2 节与下式 (5.1.2) 均采用正序。
此处据对象型别恢复。见疑误 A-0256、A-0257。}
以 $\mathcal L(\Omega(L^{(0)},L^{(1)});\ell_0)$
表示以 $\ell_0$ 为基点的环路空间；
$\gamma$ 给出的共轭拼接
\[
w\longmapsto\gamma\#w\#\gamma^{-1}
\]
诱导同胚
\[
\gamma_*:
\mathcal L(\Omega(L^{(0)},L^{(1)});\ell_0)
\longrightarrow
\mathcal L(\Omega(L^{(0)},L^{(1)});\ell_0').
\]
这里 $\gamma^{-1}$ 是 $\gamma$ 的反向道路，
即 $\gamma^{-1}(t)=\gamma(1-t)$。
直接检验可知，$\gamma_*$ 诱导同构
\begin{equation}
\gamma_*:
\pi_1(\Omega(L^{(0)},L^{(1)});\ell_0)
\longrightarrow
\pi_1(\Omega(L^{(0)},L^{(1)});\ell_0'),
\tag{5.1.2}\label{eq:5-1-2}
\end{equation}
并且保持能量与马斯洛夫指标：
\[
E_{\ell_0'}\circ\gamma_*=E_{\ell_0},
\qquad
\mu_{\ell_0'}\circ\gamma_*=\mu_{\ell_0}.
\]
因此，式~\eqref{eq:5-1-2}
又诱导同构
\[
\gamma_*:\Pi(L^{(0)},L^{(1)};\ell_0)
\longrightarrow
\Pi(L^{(0)},L^{(1)};\ell_0').
\]
综上可得：

\begin{proposition}{5.1.3}\label{prop:5-1-3}
设 $\ell_0'\in\Omega(L^{(0)},L^{(1)})$
与 $\ell_0$ 位于同一连通分支。
每选取一条连接二者的道路
$\gamma:[0,1]\to\Omega(L^{(0)},L^{(1)})$，
便确定同构
\[
\gamma_*:\Pi(L^{(0)},L^{(1)};\ell_0)
\longrightarrow
\Pi(L^{(0)},L^{(1)};\ell_0').
\]
它仅依赖于 $\gamma$ 的同伦类，且满足
\[
E_{\ell_0}=E_{\ell_0'}\circ\gamma_*,
\qquad
\mu_{\ell_0}=\mu_{\ell_0'}\circ\gamma_*.
\]
特别地，$\gamma_*$ 诱导诺维科夫环的同构
\[
\gamma_*:\Lambda(L^{(0)},L^{(1)};\ell_0)
\longrightarrow
\Lambda(L^{(0)},L^{(1)};\ell_0'),
\]
该同构同时保持滤过与分次，
并且不依赖于连接 $\ell_0$ 与 $\ell_0'$
的道路 $\gamma$ 的选取。
\end{proposition}

因此，只要 $\ell_0'$ 与 $\ell_0$
位于同一连通分支，
两个环
$\Lambda(L^{(0)},L^{(1)};\ell_0)$
与
$\Lambda(L^{(0)},L^{(1)};\ell_0')$
之间就有典范同构。

现在定义它到万有诺维科夫环的嵌入。
\begin{definition}{5.1.4}\label{def:5-1-4}
定义
\[
I_{L^{(0)},L^{(1)};\ell_0}:
\Lambda(L^{(0)},L^{(1)};\ell_0)
\longrightarrow
\Lambda_{\mathrm{nov}}[e^{1/2}]
\]
为
\[
\begin{aligned}
&I_{L^{(0)},L^{(1)};\ell_0}
\left(
\sum_{\beta\in\Pi(L^{(0)},L^{(1)};\ell_0)}
a_\beta[\beta]\right)\\
&\qquad=
\sum_{\beta\in\Pi(L^{(0)},L^{(1)};\ell_0)}
a_\beta T^{E_{\ell_0}(\beta)}
e^{\mu_{\ell_0}(\beta)/2}.
\end{aligned}
\]
\end{definition}

若 $L^{(0)}$ 与 $L^{(1)}$ 均可定向，
$\mu_{\ell_0}(\beta)$ 总是偶数，
所以 $I_{L^{(0)},L^{(1)};\ell_0}$
的像实际上落在 $\Lambda_{\mathrm{nov}}$ 内。
一般情形下，则须把万有诺维科夫环
扩张为 $\Lambda_{\mathrm{nov}}[e^{1/2}]$
才能定义该映射。
\translatornote{原书实体页 278
先把定义 5.1.4 的靶写成 $\Lambda_{\mathrm{nov}}$，
随即说明非定向情形需要 $e^{1/2}$；
译文的箭头采用适用于一般情形的扩张靶。
见疑误 A-0258。}

下面两个引理是引理 5.1.1
在拉格朗日子流形对上的对应形式，
由定义即可检验。

\begin{lemma}{5.1.5}\label{lem:5-1-5}
在上述两个诺维科夫环
$\Lambda(L^{(0)},L^{(1)};\ell_0)$
与
$\Lambda(L^{(0)},L^{(1)};\ell_0')$
之间的典范同构下，
有
\[
I_{L^{(0)},L^{(1)};\ell_0}
=I_{L^{(0)},L^{(1)};\ell_0'}.
\]
\end{lemma}

\begin{lemma}{5.1.6}\label{lem:5-1-6}
设 $\psi:M\to M'$ 是辛微分同胚，
且对 $i=0,1$ 均有
$\psi(L^{(i)})=L^{(i)\prime}$。
则存在同构
\[
\psi_*:\Pi(L^{(0)},L^{(1)};\ell_0)
\longrightarrow
\Pi(L^{(0)\prime},L^{(1)\prime};\psi(\ell_0)),
\]
进而有同构
\[
\psi_*:\Lambda(L^{(0)},L^{(1)};\ell_0)
\longrightarrow
\Lambda(L^{(0)\prime},L^{(1)\prime};\psi(\ell_0)).
\]
此外，它满足
\[
I_{L^{(0)\prime},L^{(1)\prime};\psi(\ell_0)}
\circ\psi_*
=I_{L^{(0)},L^{(1)};\ell_0}.
\]
\end{lemma}
